\section{Diskussion}
\label{sec:Diskussion}

\subsection{Invertierender Linearverstärker}
Zunächst ist zu erkennen, dass alle verwerteten Frequenzgänge den theoretischen 
Erwartungen entsprechen. Es ist grundsätzlich ein konstanter Plateaubereich 
identifizierbar. Ebenso weichen die Messwerte der Abfallgeraden kaum vom Graphen 
ab (mit Ausnahme von Messung 2, wo die letzten drei Werte Abweichungen aufweisen),
was für eine präzise Messung spricht.
\begin{table}[H]
    \centering
    \caption{Ermittelte Verstärkungen und die Abweichung zum Theoriewert.}
    \label{tab:d1}
    \sisetup{table-format=1.2}
    \begin{tblr}{
        colspec = {S[table-format=1.0] S[table-format=1.4] S[table-format=2.4] S[table-format=2.2]},
        row{1} = {guard, mode=math},
        column{2} = {halign=c},  % Zentriere Spalte 2
        column{3} = {halign=c},  % Zentriere Spalte 3
        column{4} = {halign=c},  % Zentriere Spalte 4
    }
    \toprule
    \text{Messung} & V_{exp} & V_{theo} & \text{Prozentuale Abweichung} \\
    \midrule
    1 & 1.4857 & 1.5   & 0.96 \\
    2 & 9.7143 & 10.0  & 2.94 \\
    3 & 4.9841 & 4.7   & 6.04 \\
    4 & 3.4286 & 3.133 & 9.44 \\  
    \bottomrule
    \end{tblr}
\end{table}
\noindent In \autoref{tab:d1} befindet sich der Vergleich der Werte für die 
Verstärkung. Trotz der genauen Graphen kommt es teils zu signifikannten Unterschieden 
zwischen experimentell und theoretisch ermittelten Werten. Jenes kann auf den 
Aspekt der Beschaffenheit des Operationsverstärkers zurückgeführt werden. Da es 
sich nicht um einen Idealen Verstärker handelt, ist etwa die Leerlaufverstärkung 
endlich und die Bandbreite begrenzt, was zur Folge haben kann, dass systematische 
Fehler entstehen. Möglicherweise haben nicht vollständig kompensierte Eingänge 
die Messungen in Form von Restverzerrungen beeinflusst. Messung 1 zeigt, dass
unter geeigneten Bedingungen durchaus sehr präzise Ergebnisse zu erwarten sind,
darauffolgende Messungen nehmen an Genauigkeit ab.
\vspace{0.5cm}
\\
\noindent Die Untersuchung der Phasenbeziehungen kann als erfolgreich gewertet 
werden, auch wenn die Dekandenmessungen bessere Werte hätten liefern können.
Für eine einfachere Identifikation des Graphenverlaufs wären andere 
Frequenzintervalle hilfreich gewesen. Besonders bei den ersten beiden
Durchgängen fehlt es für höhere Frequenzen an Werten für die Phase der Signale.
Trotzdem verändert sich die Phase am Ausgangsoszilloskop (ausgehend von 180°
oder -180°) um 90°. In den Plots der Messwerte ist zu erkennen, das höhere
Frequenzen zunehmend zu einer Phasenveränderung führen. Somit ergibt sich der
erwartete Verlauf eines Tiefpasses. 


\subsection{Der Umkehr-Integrator-und Differenzierer}
Für die Ausgleichsrechnung sollte sich für die Steigung im Idealfall ein Wert 
von $m=-1$ einstellen. Experimentell ergibt sich ein Wert von
\begin{equation*}
    m_{\text{exp}} = -0.9193 \pm 0.0206,
\end{equation*}
was einer Abweichung von $8.07 \%$ bzw. $6.01 \%$ entspricht. Derartige Abweichungen 
sind plausibel, da laut Datenblatt des LM741 Widerstände über eine Toleranz von 
$5 \%$ verfügen \cite{datenblatt}. Für die integrierte Sinus-, Dreieck- und 
Rechteckspannung ergeben sich ebefalls plausible Ergebnisse. Das Sinussignal 
wird um knapp 90° versetzt, das Rechtecksignal entwickelt sich zu der erwarteten 
Dreieckspannung und die Dreieckspannung wird zu einer Sinusspannung.
\vspace{0.5cm}
\\
\noindent Beim Differenzierer sind die Abweichungen größer. Eine ideale
Proportionalität wäre durch eine Steigung in der Regression von $m=1$ gegeben.
Erreicht wird 
\begin{equation*}
    m_{\text{exp}} = 0.8414 \pm 0.0428.
\end{equation*}
Die Differenzen zum Idealwert betragen $15.86 \%$ und inklusive Fehlertoleranz 
$11.58 \%$. Diese erhöhten Werte im Vergleich zum Integrator könnten aus den 
gleichen Gründen wie beim Linearverstärker kommen: unkompensierte Eingänge 
verzerren Möglicherweise die Messung. Am Ausgangsozilloskop zeigen sich die
charakteristischen Ableitungseffekte. Aus einem sinusförmigen Eingangssignal 
entsteht aufgrund der Differentiation ein cosinusförmiger Verlauf, was einer 
Phasenverschiebung von 90° entspricht. Bei der Rechteckspannung als Eingang werden
an den steilen Flanken Spannungsimpulse erzeugt, welche oszillierend abklingen.
Für die Rechteckspannung kommen ähnliche Effekte zustande. Im Idealfall würde
sich hier eine Rechteckspannung ergeben.


\subsection{Schmitt-Trigger}
Für die Hysteresebreite des Schmitt-Triggers sollte sich theoretisch ein Wert von 
\begin{equation*}
    \Delta U_{\text{theo}} = \qty{2,727}{\volt}
\end{equation*}
einstellen. Im Vergleich zum ermittelten Wert über das Oszilloskop mit
\begin{equation*}
    \Delta U_{\text{exp}} = \qty{2.736 \pm 0.002}{\volt}.
\end{equation*}
bedeutet das eine Abweichung von lediglich $0.33 \%$. Das Experiment kann daher 
als sehr gelungen bewertet werden. Die hervorragende Übereinstimmung bestätigt
die Theorie und zeigt, dass der Schmitt-Trigger im Gegensatz zu linearen
Verstärkerschaltungen oder Umkehrschaltungen sehr robust gegenüber realen
Einschränkungen ist.